\documentclass[11pt,a4paper]{article}

\usepackage[margin=1in]{geometry}
\usepackage{hyperref}
\usepackage{booktabs}
\usepackage{amsmath}
\usepackage{amssymb}
\usepackage{graphicx}
\usepackage{array}
\usepackage{longtable}
\usepackage{xcolor}
\usepackage{setspace}
\usepackage{caption}
\usepackage{subcaption}
\usepackage{natbib}
\usepackage{microtype}
\usepackage{parskip}
\usepackage{tabularx}
\usepackage{multirow}
\usepackage{float}

\hypersetup{
  colorlinks=true,
  linkcolor=blue,
  citecolor=blue,
  urlcolor=blue
}

\title{Multi-Omic Data Integration via Decentralized AI Agents\\
for the Prediction of Phenotypic Drug Resistance in Cancer}

\author{
  Shrinithi Natarajan \\
  \texttt{snatara8@jh.edu} \\
  Department of Biotechnology, \\
  Johns Hopkins University \\
  Baltimore, MD 21210, USA
  \and
  Dr.\ Bhaskarjyoti Das \\
  \texttt{bhaskarjyoti01@gmail.com} \\
  Department of AI\&ML, \\
  PES University \\
  Bangalore, India
}

\date{}

\begin{document}

\maketitle

\begin{abstract}
The prediction of phenotypic cancer drug resistance from multi-omic data remains a critical
unsolved problem in precision oncology. Existing computational frameworks either sacrifice
interpretability for predictive accuracy or lack grounded conflict-resolution mechanisms when
molecular evidence layers disagree. This study introduces a Decentralized Multi-Agent System
(DMAS) that models the deliberative structure of a Multidisciplinary Tumor Board: four specialist
LLM agents, each linked to a dedicated Model Context Protocol (MCP) server and a distinct
biological database (DepMap/CCLE, CIViC, KEGG, CTRPv2/GDSC2, ChEMBL), engage in a
two-round structured peer debate governed by a five-tier biological axiom hierarchy. Each agent
produces a structured EvidencePack with a formula-grounded confidence score (GCS) derived from
biomarker citation rate, evidence availability, modality tier, and signal strength. Conflicts are
resolved through endorsement-primary, tier-secondary lexicographic ordering.

Evaluated on 31 literature-validated gold-standard (cell line, drug) pairs with published
resistance phenotypes, the full system achieves an AUROC of 0.819, Cohen's $\kappa$ of 0.733,
AUPRC of 0.766, 100\% coverage (all 31 cases decisively resolved), 87.1\% accuracy (27/31
correct), and a sensitivity (TPR) of 1.000 with zero false negatives. Fourteen ablation conditions
isolate the contributions of the debate protocol, axiom hierarchy, MCP grounding, and each
individual agent. Claim-level faithfulness across 194 atomic claims is 86.6\%, with cross-tier
leakage identified as the primary failure mode rather than factual hallucination. The architecture
provides a scalable, auditable, and biologically faithful framework for personalized cancer therapy
decision support.
\end{abstract}

\section{Introduction}
\label{sec:intro}

For years, the clinical management of cancer was primarily determined by physiological and
anatomical origin, with focus on identifying the specific type of cancer present~\cite{hein2025,
liu2025multiomics}. With the advent of precision oncology, there has been a shift in the medical
paradigm to discover unique molecular characteristics of a tumor within a patient~\cite{hein2025}.
Multiomics, which is the integration of genomics, proteomics, transcriptomics, epigenomics and
other omics layers, is known to play a crucial role in such transitions by providing a unique
molecular signature that dictates the phenotypic drug resistance~\cite{du2024}.
Multidisciplinary Tumor Boards (MDT) or Molecular Tumor Boards (MTB) are useful in aiding
the integration of such complex and high dimensional datasets to form meaningful treatment
plans. These boards are conducted by human specialists such as oncologists, geneticists,
pathologists and pharmacologists, who debate and resolve complex and often conflicting clinical
signals~\cite{han2025,lorenzo2026}. However, as the dimensionality of such data grows
exponentially, such tumor boards face severe constraints with respect to throughput and cognitive
overload.

There are a lot of cases where targeted therapies fail to fully overcome drug resistance at
the individual patient level. This leads to deaths related to cancer to exceed over 90\% in patients
receiving targeted therapies~\cite{mao2025}. For instance, in BRAF-mutant melanoma, majority
of the patients are known to develop resistance within approximately two years of receiving BRAF
inhibitor therapy~\cite{proietti2020}. In the same way, in EGFR-driven lung cancer, the
resistance to first and second-generation EGFR inhibitors emerges routinely within 9 to 13
months~\cite{rotow2017}. With the high diversity in underlying molecular mechanisms of
resistance (target gene mutations, alternative signaling pathways and epigenetic silencing of
transcriptional programs~\cite{housman2014}), no single omic layer can successfully capture the
diversity. This requires the integration of multiple biological data layers~\cite{liu2025multiomics}.
Yet, the critical challenge still remains in resolving disagreements among these evidence layers
when conflicting signals emerge.

Existing computational methods in Cancer Drug Response (CDR) prediction are known to fall
into the following categories:
\begin{itemize}
  \item \textbf{Deep Learning and Fusion Frameworks:} Frameworks like MSDRP and DRN-CDR make
  use of deep residual networks and similarity network fusion for capturing non-linear features
  from multi-omic data~\cite{saranya2024,zhao2023}. However, such models still function as
  ``black-boxes''~\cite{almutiri2023}. Their high predictive performance does not provide a
  mechanistic reasoning or causal explanations for clinical decision support~\cite{chen2024}.

  \item \textbf{Monolithic Large Language Models (LLMs):} CellHit and other generative AI
  platforms use natural language reasoning for drug response predictions~\cite{parca2025,
  carli2025}. However, in such designs, the use of monolithic LLMs to process dense multi-omic
  payloads makes the framework vulnerable to context overwhelm and attention
  bias~\cite{zhou2025scoping,meng2024}. They fixate on the prominent DNA mutations (like BRAF
  V600E) while being blind to RNA bypass pathways. This results in inaccurate drug response
  predictions~\cite{kim2025}.

  \item \textbf{Multi-Agent Systems (MAS):} EvoMDT, HAO and BioAgents are architectures that
  design agents specific to certain roles to reduce context
  dilution~\cite{liu2026evomdt,codella2025,mehandru2025}. However, such systems with multiple
  agents lack grounded biological rules for conflict resolution. This forces the agents to resolve
  conflicts using social negotiation, weighted voting or by making use of a ``judge''
  model~\cite{kaesberg2025}. Such systems are still susceptible to sycophancy and the problem of
  ``Silent Agreement''. In these scenarios, herding effects and systemic conformity suppress the
  independent reasoning ability of each agent~\cite{cui2025,turpin2023}.
\end{itemize}

The absence of an open-source computational framework that is built on the causal ordering of
molecular evidence layers embedded as an executable logic is a primary gap in precision oncology
AI. Mere statistical agreement among the models is not sufficient for biological accuracy. This
has been the standard consensus aggregation method that is available in existing
literature~\cite{cui2025}. While Model Context Protocol (MCP) has been used in platforms like
BioinfoMCP and MCPmed to standardize semantic contracts between LLMs and structured
datasets, the usage is just as a functional tool call rather than something leveraged for explaining
biological causality~\cite{widjaja2025,flotho2026}.

To address these challenges, this study introduces a decentralized multi-agent consensus
architecture that takes inspiration from a human tumor board~\cite{zhang2025}. This design
isolates distinct biological modalities into four independent agents: Genomics, Transcriptomics,
Pathways and Pharmacology. Each of these agents act as independent MCP clients that query
databases (DepMap/CCLE, CIViC, KEGG, CTRPv2/GDSC2, ChEMBL) using schema-validated
API tool calls to prevent probabilistic hallucinations~\cite{barretina2012,iorio2016}.

The main novelty of this architecture is the transition from open-ended linguistic negotiation to
an agent-to-agent consensus protocol that is also grounded by a formal biological Hierarchy of
Truth. This is rooted in the central dogma of molecular biology where this five-tier hierarchy
(T1 Structural Primacy $\rightarrow$ T2 Transcriptional Gate $\rightarrow$ T3 Pathway Bypass
$\rightarrow$ T4 Pharmacological Prior $\rightarrow$ T5 Statistical Consensus) governs how
conflict resolution happens. Lower-tier biological axioms override the higher-tier numbers. This
ensures that an independent agent should mathematically and structurally prove a mechanistic
override (for instance, a transcriptomic sensitivity signal that is physically nullified by an active
downstream pathway bypass).

The study establishes two primary specific aims. First, to determine if this axiom-governed
multi-agent debate can predict phenotypic cancer drug resistance better than monolithic LLMs
and classic machine learning baselines, with an AUROC $\geq 0.75$ and an AUPRC $\geq 0.70$.
Second, to isolate the individual contributions of each of the architectural components using
fourteen ablation conditions (omitting axioms, debate rounds, typed MCP tools, domain isolation,
randomized tier order, standardized majority voting, and single-agent baselines for each modality).
Each prediction that is output by this system is a structured record that contains the entire debate
trace and a logged dissent file which preserves even the minority opinions when an agent's verdict
is overridden by a higher axiom tier. The design also replaces opaque statistical confidence scores
with reproducible, biologically verifiable arguments, thus introducing a scalable and auditable
framework for personalized cancer therapy.

\section{Existing Literature}
\label{sec:lit}

\subsection{MultiOmics in Precision Oncology is Now Well Established, and Why}

The field of Cancer Drug Response (CDR) prediction has evolved from static chemical-structural
analysis to sophisticated deep learning frameworks that treat drug-cell line interactions as a complex
multimodal fusion problem~\cite{saranya2024,zhao2023,chen2024}. Current research establishes a
clear consensus that single-omics data is inherently limited due to the low correlation (0.40)
between mRNA and protein levels and that integrating genomics, transcriptomics, and epigenomics
significantly improves the refined accuracy of IC50
predictions~\cite{du2024,saranya2024,chen2024,liu2023snf}. Landmark architectures like DRN-CDR
and MSDRP have advanced the field by utilizing Deep Residual Networks (ResNets) and
Similarity Network Fusion (SNF) to capture higher-order non-linear features and mitigate the
curse of dimensionality~\cite{saranya2024,zhao2023,almutiri2023,liu2023snf}. However, a
pervasive ``interpretability-performance trade-off'' remains. While these black-box models achieve
high AUC values, they frequently fail to provide the mechanistic reasoning required for clinical
adoption, leaving the causal link between molecular features and therapeutic outcomes largely
inferred~\cite{zhao2023,chen2024,jiang2025,almutiri2023}.

\subsection{Pre-trained and Instruction Tuned LLMs for MultiOmics}

Current research in generative AI for bioinformatics has transitioned from static feature
engineering to powerful, task-specific architectures like AlphaFold and Enformer, which have
revolutionized structural biology and genomic
interpretation~\cite{hein2025}. However, these ``black-box'' models often struggle with
multi-omic integration, frequently losing biological nuance when reconciling conflicting signals.
While a microservices-driven paradigm~\cite{elsayed2026} improves computational scalability and
throughput by up to 46\%, it still faces a semantic interoperability gap. Individual AI services
can exchange data but cannot autonomously debate or prioritize signals based on the specific
biological context of a tumor.

The transition from Automated Machine Learning (AutoML) to reasoning-based autonomy
represents a pivotal shift in computational biology, as Large Language Model (LLM) agents begin
to replicate the ``plan-act-reflect'' cycle of human experts~\cite{martinek2025,qu2025}. Current
research highlights that while multi-omics integration significantly outperforms single-omic models
in predictive accuracy, it remains hindered by a profound Interpretability
Gap~\cite{zack2025,vavekanand2026}. Landmark datasets like Biology-Instructions and platforms
such as PromptBio have attempted to bridge this by instruction-tuning models on millions of
biological sequences, yet these systems frequently encounter generalization collapse when faced
with the 3D complexity of genomic data or the high-dimensional sparsity of single-cell
omics~\cite{martinek2025,he2025,ali2025,yang2025llm4cell,yang2025promptbio}. Furthermore, a
persistent biological bias exists due to a lack of population diversity in training data, which leads
to reduced generalizability across different ethnic groups~\cite{zack2025,yang2025llm4cell,
vavekanand2026}.

\subsection{LLMs as Treatment Recommenders}

Recent scoping reviews and reproducibility studies~\cite{zhou2025scoping,meng2024} demonstrate
that while LLMs show human-level potential in passing medical exams, their transition to
real-world treatment recommendation is hindered by a clinical utility gap. Specifically, current
models struggle with the iterative and messy nature of longitudinal Electronic Health
Records (EHRs)~\cite{zhou2025scoping,lin2025} and often exhibit an over-caution bias, as seen
in large-scale studies where LLMs significantly lagged behind physicians in diagnostic
specificity~\cite{lin2025,williams2024}. Furthermore, the hallucination and safety bottleneck
remains a primary concern. Even advanced models often favor probabilistic patterns over deep
disease pathophysiology~\cite{meng2024}, leading to recommendations that may be linguistically
fluent but biologically inaccurate. While hybrid neuro-symbolic
architectures~\cite{galitsky2024} and Retrieval-Augmented
Generation (RAG)~\cite{zhou2025scoping,lin2025} have been proposed to mitigate these risks,
most existing frameworks remain focused on clinical text summarization --- reading doctor's notes
rather than calculating the underlying molecular drivers of a
disease~\cite{lin2025,williams2024}.

\subsection{LLMs in Tumor Drug Response}

The challenge of tumor drug resistance is fundamentally a problem of information lag and
biological complexity that traditional statistical methods struggle to
resolve~\cite{mao2025}. Current research has transitioned from single-modality analysis to
multimodal AI models that integrate genomics, transcriptomics, and the tumor microenvironment
to achieve high predictive accuracy~\cite{mao2025,liu2025multiomics,wu2026,su2025}. Landmark
frameworks such as CellHit and GeneRxGPT have demonstrated that grounding Large Language
Models (LLMs) in mechanistic pathway-aware features or real-time literature retrieval
significantly outperforms static, black-box models~\cite{carli2025,lai2025}. However, despite
these advances, a significant translational gap remains. Models trained on in vitro cell lines often
fail in real-world patients due to intra-tumor heterogeneity and the Black Box problem, where
clinicians cannot decipher the underlying biological rationale behind a resistance
prediction~\cite{mao2025,carli2025,su2025}.

\subsection{LLM in MDT --- Single and Multi-Agent Systems}

The digital transformation of Multidisciplinary Tumor Boards (MDTs) has transitioned from
monolithic, rule-based clinical decision support toward decentralized, multi-agent architectures
that mirror the collaborative nature of human medical
teams~\cite{liu2026evomdt,zhang2025,han2025}. Recent literature emphasizes that
role-differentiated agents, specializing in diagnostics, treatment, and safety, significantly
outperform single-model approaches by reducing context dilution and enhancing reasoning
accountability~\cite{liu2026evomdt,zhang2025,codella2025}. Frameworks such as EvoMDT and
HAO have demonstrated high factual recall and guideline concordance by utilizing specialized
agents grounded in Retrieval-Augmented Generation
(RAG)~\cite{liu2026evomdt,codella2025}. However, despite these architectural gains, the field
faces persistent bottlenecks of reasoning. Current benchmarks like MTBBench show that models
perform near-randomly on complex longitudinal tasks like recurrence forecasting, and most systems
struggle to move beyond surface-level text summarization to resolve the conflicting multi-omic
signals inherent in precision oncology~\cite{vasilev2025,codella2025,lorenzo2026}.

\subsection{Debate and Consensus Mechanisms in MAS}

The integration of Large Language Models (LLMs) into bioinformatics has shifted the field from
static algorithmic tools toward dynamic, reasoning-based multi-agent systems (MAS) capable of
managing complex, multi-step workflows~\cite{mehandru2025,yang2025rise}. As identified in
recent frameworks like BioAgents, while single-agent systems excel at conceptual genomics and
tool selection, they frequently struggle with the functional execution of code and the procedural
knowledge gap required for end-to-end analysis~\cite{mehandru2025}. Furthermore, the iterative
refinement process often touted in multi-agent collaboration has shown diminishing returns.
Repeated self-correction rounds can lead to knowledge inertia or even degrade output quality,
suggesting that linguistic debate is not a substitute for grounding in high-quality domain
data~\cite{mehandru2025,zhou2025streamline}.

This limitation is closely tied to the architectural tension between collaborative consensus and
voting paradigms. While traditional Multi-Agent Debate (MAD) frameworks treat these protocols
as static variables, recent empirical evaluations demonstrate a sharp performance dichotomy:
consensus protocols show a performance improvement of 2.8\% on factual, knowledge-dense tasks,
whereas voting protocols yield a 13.2\% optimization on abstract reasoning
tasks~\cite{kaesberg2025}. However, scaling conversational rounds blindly deteriorates system
accuracy due to ``problem drift'' and ``sycophantic convergence''~\cite{kaesberg2025}, alongside
an inherent unreliability tax where multi-agent frameworks consume 15 to 50 times more tokens
than standalone models~\cite{spieser2026}. To counteract this computational overhead,
contemporary architectures utilize heterogeneous agent groups to introduce epistemic diversity and
combat shared bias~\cite{liu2026heterogeneous}. They also deploy sparse, progressive
communication topologies such as localized intra-group debate paired with hierarchical
inter-group synthesis which can reduce total token expenditure by up to 46.9\% without
sacrificing logical reasoning depth~\cite{liu2024groupdebate}.

Crucially, standard consensus-driven aggregation models rely heavily on a final ``judge'' model or
majority voting, working under the flawed assumption that agreement equates to biological
accuracy~\cite{cui2025}. This architecture is highly susceptible to the ``Silent Agreement''
problem, where systemic conformity and herding effects suppress critical, independent reasoning
paths~\cite{cui2025}. When no clear consensus exists, standard majority voting becomes
statistically random and inequitable~\cite{cui2025}. While advanced frameworks have introduced
mathematically rigorous aggregation tools, such as using Kendall's coefficient of concordance (W)
to identify discordant agents or applying multi-agent reinforcement learning (PPO) to optimize
consensus paths, these methods frequently experience latency spikes and require up to 252 GB of
memory~\cite{han2025}.

This technical barrier separates the current lab-to-live execution gap, where systems performing
at a 90\% accuracy tier on static examinations experience a catastrophic drop to 45--69\% when
subjected to practice-based, unstructured clinical data~\cite{spieser2026}. Furthermore, existing
autonomous workflows remain largely reactive and restricted by a top-down supervisor-worker
topology, such as PromptBio's supervisor routing engine, which exhibits severe performance
degradation when generating code for complex, multi-step survival
analyses~\cite{yang2025promptbio}. This underscores the critical need for decentralized,
peer-to-peer agent interaction spaces that implement native, reasoning-heavy backbones alongside
standardized interfaces to safely bridge data operations with true biological causality.

\subsection{MCPs in Bioinformatics and MAD}

The transition from human-mediated command-line tools to autonomous, agent-ready
infrastructures represents a critical frontier in computational
biology~\cite{widjaja2025,flotho2026}. As identified in recent frameworks like BioinfoMCP, the
bottleneck in AI-driven science has shifted from raw reasoning capability to the manual
engineering required to interface Large Language Models (LLMs) with legacy bioinformatics
tools~\cite{widjaja2025}. While the Model Context Protocol (MCP) provides a standardized
semantic contract that moves beyond basic APIs to enable machine-actionable data stewardship,
current implementations often struggle with Hardware-Software Co-optimization and a lack of
robust Error-Aware Recovery for biological edge cases~\cite{widjaja2025,flotho2026}.
Furthermore, a significant research gap remains in providing an audit trail for high-stakes clinical
scenarios, as most current MCP-enabled agents focus on functional execution rather than the
mechanistic explainability required for regulatory trust~\cite{flotho2026}.

The shift from single-model prompting toward collaborative multi-agent architectures has
highlighted the potential of Multi-Agent Debate (MAD) to mitigate hallucinations and enhance
reasoning through iterative peer review. However, as noted in the M3MAD-Bench
study~\cite{li2026}, current frameworks often suffer from divergent noise in unimodal tasks and a
plateau effect, where simply increasing the rounds of debate fails to reach a more accurate
convergence. A critical research gap remains in the reliance on fixed roles and rigid protocols
that fail to capture environment-driven feedback, leading to ``collective delusion'' where agents
mutually reinforce incorrect assumptions due to a lack of grounding in external, high-fidelity data.

The preceding review identifies a critical divergence in the field of precision oncology. While
predictive accuracy in multi-omic integration is rising, clinical interpretability and causal
grounding remain stagnant. Current architectures, ranging from monolithic LLMs to RAG-based
clinical assistants, consistently struggle with probabilistic hallucination, context dilution, and a
reliance on surface-level text rather than raw biological data. To bridge this gap, this thesis
introduces a Decentralized Multi-Agent Consensus Framework. By utilizing the Model Context
Protocol (MCP), the framework creates a live semantic bridge to structured databases
(DepMap/CCLE, CIViC, KEGG, CTRPv2/GDSC2, ChEMBL), enabling agents to bypass uncurated
text and reason directly over raw numerical biomarkers like Z-scores and mutation frequencies.

The core innovation of this work lies in the transition from open-ended linguistic debate to a
Hierarchy of Truth protocol. Unlike traditional multi-agent systems that rely on weighted voting
or LLM-as-a-judge mechanisms, this framework implements a deterministic consensus logic
grounded in the Central Dogma of molecular biology. By requiring agents to prove mechanistic
overrides, such as a DNA structural binding failure physically nullifying a transcriptomic
sensitivity signal, the system replicates the deliberative rigor of a Multidisciplinary Tumor
Board (MDT). This approach directly solves the ``black-box'' problem and the ``cold-start''
challenge for new drugs by grounding every clinical verdict in verifiable biological axioms rather
than pattern-matching. In doing so, this research provides a scalable, auditable, and biologically
faithful architecture for the next generation of personalized cancer therapy.

\section{Methodology}
\label{sec:methodology}

This study introduces a Decentralized Multi-Agent System (DMAS) which has four specialized LLM
agents. Each agent operates on a distinct omic modality and independently analyzes a (cell line,
drug) pair. They then submit a structured evidence pack which is used to engage in a two-round
peer debate. The outcome of the debate is governed by an explicit biological axiom hierarchy.
This architecture is grounded in today's clinical practice in the form of Multidisciplinary Tumor
Boards (MDTs) where oncologists, geneticists, pathologists, pharmacists, etc.\ debate and resolve
conflicting clinical signals~\cite{han2025,lorenzo2026}. The central dogma of molecular biology
which establishes that DNA-level structural alterations constitute the highest-level evidence of
drug-target interaction also plays a crucial role here~\cite{han2025,liu2025multiomics}.

This design is implemented using Python programming language using FastMCP for a structured
tool access. Google Vertex AI (specifically Gemini 3.1 Flash Lite) is used for the LLM backend.
SQLite databases are built using four publicly available bioinformatics resources. Every data
access by the agents is mediated using Model Context Protocol (MCP) tool calls. This is
essential for converting hallucinations from a probabilistic failure mode to a detectable protocol
violation, that is, an agent cannot assert that a mutation exists if the genomics server does not
return that data~\cite{flotho2026,widjaja2025}.

% TODO: INSERT FIGURE --- Figure 1: Five-layer framework architecture. Data flows from public databases (Layer 1) through MCP-validated tool APIs (Layer 2) to specialist LLM agents (Layer 3), which engage in a two-round debate protocol (Layer 4), producing a structured output with a full audit trail (Layer 5).
\begin{figure}[H]
  \centering
  \caption{Five-layer framework architecture.}
  \label{fig:architecture}
\end{figure}

\subsection{Dataset Sources}

The architecture uses four publicly available bioinformatics databases selected to provide orthogonal evidence for the prediction of drug sensitivity. It is crucial to use multiple data sources as single-omics data is heavily limited owing to its low correlation (approximately 0.40) between mRNA and protein levels. Any single data modality only captures a single dimension of the highly complex cellular state that determines drug response~\cite{saranya2024,zhao2023}.

\subsubsection{Genomic Data --- CCLE/DepMap Somatic Mutations and Copy Number Variation}

The Cancer Cell Line Encyclopedia (CCLE) / DepMap project is used to retrieve somatic mutation profiles and copy number variations (CNVs). This database is maintained by the Broad Institute and the Wellcome Sanger Institute~\cite{barretina2012}. CCLE provides whole-exome sequencing mutation calls and CNV segment data for around 1,800 cancer cell lines which spans more than 30 cancer lineages. The COSMIC Cancer Gene Census annotates the driver mutations. For every mutation, the cell line identifier, protein change (HGVS notation), the gene symbol, mutation type (missense, frameshift, nonsense, splice site), CIViC description and the driver status are retained.

\textbf{CIViC enrichment.} Mutation records are enriched using drug-specific clinical evidence at the query time using Clinical Interpretation of Variants in Cancer (CIViC) database~\cite{griffith2017}. It provides a structured evidence statement for every (gene, variant, drug) triplet, linking it to either resistant or sensitive phenotypes with evidence-level ratings from A to E. Fuzzy drug-name normalization is used to perform the enrichment by reducing the query drug and the CIViC drug field to uppercase alphanumeric strings after which a bidirectional containment check is applied. For instance, it would allow Vemurafenib to match CIViC entries for PLX4032. This type of evidence is the highest priority signal within the T1 tier where a direct clinical evidence that links the response of a drug to a mutation supersedes all other genomic inferences.

\subsubsection{Transcriptomic Data --- CCLE mRNA Expression}

The DepMap public portal~\cite{barretina2012} is used to retrieve CCLE RNA-sequencing profiles. Expression values are recorded as $\log_2(\text{TPM}+1)$ that are also z-score normalized for every gene across all the cell lines. This produces a cell-line-specific deviation score. The transcriptional silencing threshold is defined by a z-score value which is below $-2.0$. So a gene with $z < -2.0$ is considered to be functionally absent from the given cell line's transcriptional programme. This constitutes a resistance mechanism for drugs that require target expression for efficacy. Similarly, a z-score above $+1.5$ is considered as a signal for overexpression for drugs that are expression-driven. High-expression gene lists are pre-filtered at $z \geq 0.5$ for scoring an active pathway and at $z \geq 1.0$ for a bypass gene detection.

\subsubsection{Transcriptional Silencing Threshold}

The T2 agent applies a silencing gate before reasoning over expression data. The threshold is defined as:
\begin{equation}
\text{silencing}(\text{gene}, \text{cell line}) =
\begin{cases}
1, & z_{\text{gene}}(\text{cell line}) < \theta, \\
0, & \text{otherwise},
\end{cases}
\quad \theta = -2.0 \;\text{(SILENCING\_THRESHOLD)}.
\end{equation}
If $\text{silencing}(\text{target gene}, \text{cell line}) = 1$, the verdict is forced to RESISTANT, regardless of the drug mechanism type.

\subsubsection{Pathway Data --- KEGG Signalling Pathways}

The Kyoto Encyclopedia of Genes and Genomes (KEGG)~\cite{kanehisa2000} is used to retrieve signalling pathway gene membership and topology data. For every pathway, the database stores the complete gene membership list and any intra-pathway roles (such as kinase, effector, transcription factor, etc). Bypass routes --- alternative signalling genes within the same pathway that can maintain downstream oncogenic signalling when a drug target is blocked --- are stored as directed (pathway id, blocked gene) $\rightarrow$ bypass gene triplets. This database covers 27 cancer-relevant signalling pathways including MAPK/ERK, PI3K/mTOR, EGFR, ABL and regulation of apoptosis.

\subsubsection{Pharmacological Data --- CTRPv2 / GDSC2 IC50 Profiles}

Two large-scale pharmacogenomic screens are used to extract historical drug sensitivity data. The Cancer Therapeutics Response Portal v2 (CTRPv2) measures drug sensitivity across around 860 cancer cell lines and 481 compounds~\cite{rees2016}. GDSC2 (Genomics of Drug Sensitivity in Cancer, release 2) provides IC50 profiles across 809 cell lines and 198 compounds, measured at the Wellcome Sanger Institute~\cite{iorio2016}. The retrieved IC50 values are converted into z-scores per drug across all the tested cell lines where $z < -0.5$ indicates a SENSITIVE response and $z > +0.5$ indicates a RESISTANT response and $|z| \leq 0.5$ is labelled as UNCERTAIN. ChEMBL is used to supplement the data for drug target information and mechanism of action~\cite{mendez2019}.

% TODO: INSERT FIGURE --- Figure 2: Four external data sources underpinning the framework.
\begin{figure}[H]
  \centering
  \caption{Four external data sources underpinning the framework. Each is publicly available for academic use and covers a distinct biological evidence layer.}
  \label{fig:datasources}
\end{figure}

\subsection{Data Preprocessing and Database Construction}

All four data sources are preprocessed and stored as local SQLite datasets to enable fast, offline tool calls without depending on external APIs at inference time. This ensures reproducibility and cost control.

\subsubsection{Genomics Database (\texttt{genomics.db})}

CCLE's raw mutation calls are filtered to retain only coding mutations with variant allele frequency $\geq 0.1$. Driver status is derived from the COSMIC Cancer Gene Census v99. Protein changes are normalized by stripping the \texttt{p.} prefix and uppercasing (e.g.\ \texttt{p.V600E} $\rightarrow$ \texttt{V600E}) for fuzzy matching with CIViC variant nomenclature. CNV segment files are processed to gene-level calls: amplification (segment mean $> 2.0$), hemizygous deletion (0.5--1.0), and homozygous deletion ($< 0.5$). CIViC predictive evidence is stored in a separate \texttt{predictive\_evidence} table and enriched at query time to enable drug-specific filtering without pre-joining all possible triples.

\subsubsection{Transcriptomics Database (\texttt{transcriptomics.db})}

CCLE RNA-sequencing data ($\log_2(\text{TPM}+1)$) is z-score normalized per gene across all 1,800+ cancer cell lines. Z-scores and TPM values are stored in a flat expression table indexed on (cell line, gene). A materialized view of highly expressed genes ($z > 0.5$) is maintained for fast pathway queries. The silencing threshold $z = -2.0$ is stored as a module-level constant.

\subsubsection{Pathway Database (\texttt{pathways.db})}

KEGG REST API data is stored in three tables: \texttt{pathway\_meta} (identifier, name, category), \texttt{pathway\_genes} (membership with role annotations), and \texttt{bypass\_routes} ((pathway id, blocked gene) $\rightarrow$ bypass gene triplets). A pathway activity score (fraction of pathway member genes with $z \geq 0.5$) is computed at inference time. Pathways with activity score below 0.4 are excluded from bypass analysis.

\subsubsection{Pharmacology Database (\texttt{pharmacology.db})}

IC50 values from CTRPv2 and GDSC2 are converted to natural-log scale and z-score normalised per drug:
\begin{equation}
z_{\text{drug}}(\text{cell line}) = \frac{\ln(\text{IC50}) - \mu_{\text{drug}}}{\sigma_{\text{drug}}}, \quad
\text{label}(z) =
\begin{cases}
\text{SENSITIVE}, & z < -0.5, \\
\text{RESISTANT}, & z > 0.5, \\
\text{UNCERTAIN}, & |z| \leq 0.5.
\end{cases}
\end{equation}
Pre-computed sensitivity labels are stored alongside raw z-scores. ChEMBL drug target information and mechanism of action annotations are stored in a \texttt{drug\_info} table. A separate LLM response cache keyed on (model, prompt hash) eliminates redundant LLM calls during iterative evaluation runs.

\subsubsection{Case Construction and Label Derivation}

We evaluate the system on 31 literature-validated gold-standard (cell line, drug) pairs. Each case was selected on the basis that its resistance or sensitivity phenotype has been established in at least one peer-reviewed publication. Unlike IC50-derived labels from pharmacogenomic screens, which can conflate cytostatic growth arrest with genuine resistance, these ground-truth labels reflect experimentally confirmed mechanistic outcomes. The 31 cases span diverse cancer lineages (breast, lung, colorectal, melanoma, lymphoma, and others) and drug classes (CDK4/6 inhibitors, EGFR inhibitors, BRAF inhibitors, PARP inhibitors, and others), and include both SENSITIVE ($n=18$) and RESISTANT ($n=13$) phenotypes.

\subsection{Framework Architecture}

The framework is organised into five hierarchical layers, each with a unique responsibility. The key architectural principle is the separation of evidence access from evidence reasoning: agents never query databases directly. All data access is managed through MCP tool calls, which return JSON-typed objects and are logged, enabling post-hoc hallucination detection.

\begin{table}[H]
\centering
\caption{Hierarchical layers of the proposed framework.}
\label{tab:layers}
\begin{tabularx}{\textwidth}{llX}
\toprule
\textbf{Layer} & \textbf{Component} & \textbf{Responsibility} \\
\midrule
1: Data Substrate & 4$\times$ SQLite databases & Store pre-processed public omics data; serve as ground truth for hallucination detection. \\
2: MCP Servers & 4$\times$ FastMCP instances & Expose typed tool APIs, validate inputs, return schema-conformant JSON. \\
3: Specialist Agents & 4$\times$ LLM agents & Fetch modality-specific evidence, reason over it, produce structured EvidencePacks. \\
4: Debate Protocol & DebateEngine + AxiomResolver & Coordinate two-round peer debate, resolve conflicts via endorsement-weighted axiom logic. \\
5: Evaluation Pipeline & Metrics + RunLogger + Ablation runner & Compute AUROC/AUPRC/$\kappa$/$\rho$, log the full audit trail, run ablation variants. \\
\bottomrule
\end{tabularx}
\end{table}

The Orchestrator receives a (cell line, drug) pair as input, initiates all four agent \texttt{analyze()} calls concurrently using \texttt{asyncio.gather()}, collects the resulting EvidencePacks, and passes them to the DebateEngine. Parallel agent execution is possible because agents have no shared state and carry out independent MCP tool calls. A shared LLM client enforces a rate limiter to prevent API quota exhaustion. After the debate protocol completes, the Orchestrator logs the ConsensusResult to the RunLogger.

\subsection{Model Context Protocol Servers}

The Model Context Protocol (MCP)~\cite{flotho2026,widjaja2025} is a recently standardised protocol for connecting LLMs to custom tools through typed, schema-validated function calls. Each MCP tool server is implemented using FastMCP. The design inherently enforces an anti-hallucination property: every data claim an agent makes must originate from an MCP tool call, and tool inputs and outputs are logged by the RunLogger, enabling per-case hallucination rate computation.

In conventional LLM prompting, an agent can probabilistically assert that a mutation exists or a gene is overexpressed without any verifiable connection to actual data. The MCP architecture converts this probabilistic failure mode into a detectable protocol violation: an agent cannot assert that BRAF V600E is present in a cell line unless the Genomics MCP server returned it. The guarantee is architectural, not probabilistic --- the same design principle used in database transactions.

\subsubsection{Genomics MCP Server}

Exposes three tools: \texttt{get\_mutations(cell\_line, gene, drug)}, returning all somatic mutations enriched with CIViC predictive evidence specific to the queried drug; \texttt{get\_cnv(cell\_line, gene)}, returning copy number variation calls; and \texttt{check\_mutation\_impact(gene, mutation)}, returning driver status and CIViC descriptions for a specific variant.

\subsubsection{Transcriptomics MCP Server}

Exposes three tools: \texttt{get\_expression(cell\_line, gene)}, returning TPM and z-score; \texttt{get\_high\_expression\_genes(cell\_line, z\_threshold)}, returning all genes above the threshold; and \texttt{check\_silencing(cell\_line, gene)}, returning a boolean silencing flag. The silencing threshold ($z = -2.0$) is injected from \texttt{axiom\_rules.py} to ensure consistency with the T2 agent's system prompt.

\subsubsection{Pathway MCP Server}

Exposes four tools: \texttt{find\_pathways\_for\_gene(gene)}, \texttt{get\_pathway\_genes(pathway\_id)}, \texttt{check\_bypass(pathway\_id, blocked\_gene)}, and \texttt{get\_upstream\_regulators(gene)}. The bypass check is the critical T3 evidence call: it verifies whether an alternative signalling route exists that could maintain oncogenic signalling after the drug target is blocked.

\subsubsection{Pharmacology MCP Server}

Exposes three tools: \texttt{get\_ic50(cell\_line, drug)}, returning the ln IC50, z-score, and pre-computed sensitivity label; \texttt{get\_drug\_info(drug)}, returning target genes, mechanism class, and pathway; and \texttt{get\_sensitivity\_profile(drug)}, returning population-level sensitivity distributions per mutation class.

\subsection{Specialist LLM Agents}

Each specialist agent is extended as a subclass from the BaseAgent class, which provides shared infrastructure for MCP tool calls, LLM prompt construction, response parsing, GCS computation, and RunLogger integration. All agents use the same LLM backend (Gemini 3.1 Flash Lite via Google Vertex AI), ensuring that performance differences between agents reflect differences in data modality rather than model capability.

\subsubsection{T1: Genomics Agent}

\textbf{Data Modality:} Somatic mutations, copy number variation, CIViC clinical evidence.

\textbf{Decision logic:} (1) Check for CIViC drug-specific annotation --- if present, it is decisive. (2) If driver mutation present without CIViC: UNCERTAIN. (3) CNV amplification of target oncogene: SENSITIVE signal. (4) Homozygous deletion: RESISTANT. (5) Hemizygous deletion: UNCERTAIN. (6) Wild-type in DB: UNCERTAIN with caveat that chromosomal fusions are not captured.

\textbf{Key constraint:} Wild-type status never defaults to RESISTANT.

\subsubsection{T2: Transcriptomics Agent}

\textbf{Data Modality:} mRNA expression z-scores, high-expression gene context.

\textbf{Decision logic:} Classifies the drug's mechanism type: (A) Overexpression-driven --- target expression level IS the signal, (B) Mutation-driven --- downstream pathway activity is assessed instead, (C) Mechanism-agnostic --- transcriptomics is not informative: UNCERTAIN.

\textbf{Silencing gate:} Target gene $z < -2.0$ implies functionally absent: RESISTANT signal regardless of mechanism type.

\subsubsection{T3: Pathway Agent}

\textbf{Data Modality:} KEGG signalling pathway topology, bypass route detection.

\textbf{Decision logic:} (1) Find pathways containing the drug target gene. (2) Apply pathway activity gate: skip pathways where $<40\%$ of member genes are expressed ($z \geq 0.5$). (3) For active pathways, check for bypass routes whose genes are upregulated ($z \geq 1.0$). (4) Bypass found + expressed: RESISTANT. (5) Target in pathway, no bypass: SENSITIVE signal. (6) Target not in any pathway: UNCERTAIN.

\textbf{Self-attestation gate:} Agent must answer 4 binary questions before committing to a decisive verdict. Score $< 3/4$: verdict demoted to UNCERTAIN.

\subsubsection{T4: Pharmacology Agent}

\textbf{Data Modality:} Historical IC50 z-scores, drug target and mechanism data.

\textbf{Decision logic:} Anchors verdict directly to the pre-computed IC50 label (SENSITIVE/RESISTANT/UNCERTAIN). The agent reports it faithfully and notes that higher-tier evidence (T1--T3) can override it. T4 is the most reliable signal for mechanism-agnostic drugs where genomic or pathway-level predictors are absent.

All four agents run simultaneously and produce independent EvidencePacks before any inter-agent communication occurs. This enforces the independence requirement: an agent's R1 verdict must reflect its own evidence~\cite{turpin2023,meng2024}.

\subsection{Evidence Pack Schema and Grounded Confidence Score}

\subsubsection{EvidencePack Schema}

Every agent's output is a structured EvidencePack --- a validated data class that contains the formal unit of inter-agent communication.

\begin{table}[H]
\centering
\caption{EvidencePack schema fields.}
\label{tab:evidencepack}
\begin{tabularx}{\textwidth}{llX}
\toprule
\textbf{Field} & \textbf{Type} & \textbf{Description} \\
\midrule
\texttt{agent\_id} & str & Unique agent identifier. \\
\texttt{verdict} & Enum & SENSITIVE $|$ RESISTANT $|$ UNCERTAIN. \\
\texttt{confidence} & float $[0,1]$ & Grounded Confidence Score (GCS). \\
\texttt{evidence\_tier} & Enum & T1\_STRUCTURAL $|$ T2\_TRANSCRIPTIONAL $|$ T3\_PATHWAY $|$ T4\_PHARMACOLOGICAL $|$ T5\_STATISTICAL. \\
\texttt{key\_findings} & list[Finding] & Structured cited biomarkers with source, value, and interpretation. \\
\texttt{reasoning} & str & Step-by-step natural-language reasoning chain. \\
\texttt{caveats} & list[str] & Identified limitations or missing data. \\
\texttt{conflict\_flags} & list[str] & Signals that contradict the stated verdict. \\
\texttt{peer\_scores} & dict & Populated in Round 2: mechanistic relevance scores given to peers. \\
\texttt{self\_attestation} & dict $|$ None & T3 only: four-question binary checklist and integer score. \\
\texttt{data\_status} & str $|$ None & \texttt{data\_found} $|$ \texttt{target\_wild\_type} $|$ \texttt{signal\_below\_threshold} $|$ \texttt{cell\_line\_missing}. \\
\bottomrule
\end{tabularx}
\end{table}

\subsubsection{Grounded Confidence Score (GCS)}

The GCS replaces the LLM's self-reported confidence, which has been shown to be poorly calibrated in zero-shot settings~\cite{turpin2023}, with a formula-grounded score:
\begin{equation}
\text{GCS} = 0.40H + 0.30E + 0.20T + 0.10S
\end{equation}
where $H = |\{f \in F : \text{biomarker}(f) \in R_{\text{MCP}}\}| / |F|$ (biomarker citation grounding), $E \in \{0,1\}$ (evidence availability), $T \in \{1.0, 0.8, 0.6, 0.4, 0.2\}$ for tiers T1--T5, and $S \in [0,1]$ (signal strength). For the pharmacology agent, $S = \min(|z|/3.0, 1.0)$. UNCERTAIN verdicts are capped: $\text{GCS}_{\text{final}} = \min(\text{GCS}, 0.50)$.

\subsection{Debate Protocol}

The debate protocol is modelled on clinical tumor board practice~\cite{han2025,lorenzo2026} and comprises two rounds of debate with the resolver as last resort. Unstructured LLM debates have two primary failure modes: sycophantic convergence where agents revise verdicts toward the majority without new evidence~\cite{turpin2023}, and context overwhelm where a single agent given all omics data simultaneously fails to weigh conflicting signals~\cite{liu2026evomdt}.

\subsubsection{T3 Self-Attestation Gate}

Before the T3 pathway agent's verdict is counted as decisive in the Round 1 consensus check, it must pass a self-attestation gate:
$\text{attestation score} = \sum_{k=1}^{4} a_k, \; a_k \in \{0, 1\}.$
Four questions: (1) bypass gene named and specific, (2) pathway active in the queried cell line, (3) drug mechanism relevant to pathway, (4) bypass gene distinct from primary drug target. Gate: score $< 3 \Rightarrow$ T3 verdict demoted to UNCERTAIN.

% TODO: INSERT FIGURE --- Figure 3: Debate protocol flow diagram.
\begin{figure}[H]
  \centering
  \caption{Debate protocol flow. Round 1 independent analysis feeds into a consensus check. If no consensus, Round 2 peer critique fires. If still no consensus, the AxiomResolver applies endorsement-primary, tier-secondary ordering.}
  \label{fig:debate}
\end{figure}

\subsubsection{Round 1: Independent Analysis and Consensus Check}

In the first round, all four agents run concurrently and independently. Each queries its assigned MCP servers, applies modality-specific reasoning, and produces an EvidencePack with a verdict and GCS score. After all EvidencePacks are returned, the \texttt{check\_consensus()} function determines whether consensus can be established:
\begin{itemize}
  \item \textbf{Decisive threshold:} An agent is decisive if its verdict is not UNCERTAIN. A single decisive agent is sufficient for R1 consensus.
  \item \textbf{T3 self-attestation gate:} Pathway agent verdict is decisive only if attestation score $\geq 3/4$.
  \item \textbf{T1 veto:} If the T1 agent has a decisive verdict contradicting the majority, consensus fails and the case escalates to Round 2.
  \item \textbf{Majority requirement:} Multiple decisive agents must produce a strict majority ($n > \text{len(decisive)}/2$). A 1-vs-1 split escalates to Round 2.
\end{itemize}

\subsubsection{Round 2: Peer Critique}
\label{subsec:r2}

If there is no consensus in R1, the DebateEngine invokes peer critique. In Round 2, each agent receives the R1 verdict and reasoning of every other agent. Rather than a tier-specific binary checklist, agents assign a single \textbf{mechanistic relevance score} $p \in [0.0, 1.0]$ to each peer's argument: 1.0 means the peer explicitly names a specific gene or biomarker value and explains the mechanistic link to drug response; 0.5 means relevant but only indirect mechanistic connection; 0.0 means generic reasoning with no specific connection to the case. This continuous scoring replaces an earlier binary checklist design because mechanistic relevance is inherently graded.

While agents enter Round 2 with their Round 1 verdict, three mechanisms allow revision:
\begin{enumerate}
  \item \textbf{AxiomChallenge concession:} When a higher-tier agent challenges a lower-tier agent's verdict and the lower-tier agent cannot rebut with modality-specific evidence, the conceding agent's verdict reverts to UNCERTAIN and its confidence is multiplied by 0.65.
  \item \textbf{Explicit contradiction revision:} An agent that cites a cross-modal contradiction in its R2 reasoning may revise its verdict; this revision is penalised by a factor of 0.78 on its confidence to discourage opportunistic flipping.
  \item \textbf{Peer-pressure revision:} An agent whose mean received peer score falls below 0.35 is automatically revised to the peer majority verdict at $0.70\times$ its original confidence.
\end{enumerate}

Endorsement scores aggregate the mechanistic relevance floats:
$\text{endorsement}(p_i) = \frac{1}{|A|-1} \sum_{j \neq i} p_{ji},$
where $p_{ji}$ is the score agent $j$ assigns to agent $i$'s argument.

\subsection{Axiom Hierarchy and Conflict Resolution}

The axiom hierarchy maps the causal ordering of biological evidence types into the conflict resolution mechanism. DNA-level structural alterations are the highest-fidelity determinants of drug-target interaction state, because they directly modify the physical binding site or gene expression capacity of the drug target~\cite{han2025,liu2025multiomics}. mRNA expression provides a functional-state signal but is subject to post-transcriptional regulation. Pathway topology provides mechanistic context but requires inferred connections. Pharmacological history is a population-level prior that integrates all upstream molecular factors but lacks specificity for the individual cell line.

% TODO: INSERT FIGURE --- Figure 4: The five-tier axiom hierarchy.
\begin{figure}[H]
  \centering
  \caption{The five-tier axiom hierarchy. Tiers reflect causal proximity to the drug-target interaction state.}
  \label{fig:hierarchy}
\end{figure}

\subsubsection{Resolver Algorithm}

If neither R1 nor R2 consensus is reached, the AxiomResolver selects the winning agent using a three-key lexicographic sort:
$p^* = \operatorname*{lex\,arg\,max}_{p \in C} \bigl(\text{endorsement}(p),\; \text{axiom\_weight}(\text{tier}(p)),\; \text{GCS}(p)\bigr),$
where axiom\_weight is 5, 4, 3, 2, 1 for T1--T5 respectively. The sort is lexicographic: endorsement first, then tier, then GCS. A T4 agent with very high peer endorsement can override a T1 agent with low endorsement, which is the expected behaviour --- factual grounding as judged by peers outweighs nominal tier priority.

\textbf{WEAK T1 OVERRIDE GATE:} When all of the following hold: (1) a decisive T1 agent exists, (2) GCS$(p^*) < \tau = 0.80$ (T1\_OVERRIDE\_CONFIDENCE\_THRESHOLD), and (3) all lower-tier decisive agents unanimously support a different verdict --- the confidence-weighted fallback selects $v^* = \arg\max_v \sum_{\{p: \text{verdict}(p)=v\}} \text{GCS}(p).$

\textbf{Confidence adjustment:} For each agent $a_i$, after the resolver selects winner $w$: non-winning agents have confidence multiplied by $(1-\delta)$, $\delta = 0.15$ (CONFIDENCE\_DECAY\_PER\_FLIP). Final confidence is the average of adjusted GCS scores across all agents.

\subsection{Evaluation Design and Metrics}

\subsubsection{Primary Metrics}

The framework is evaluated using four complementary metrics. UNCERTAIN predictions are excluded from AUROC/AUPRC/Spearman computation; they are included in Cohen's $\kappa$ as a third label class.

\begin{table}[H]
\centering
\caption{Primary evaluation metrics and their rationale.}
\label{tab:metrics}
\begin{tabularx}{\textwidth}{llX}
\toprule
\textbf{Metric} & \textbf{Formula} & \textbf{Rationale} \\
\midrule
AUROC & AUC of ROC curve using confidence scores and binary labels. & Measures ranking quality; threshold-independent; standard pharmacogenomics benchmark. \\
AUPRC & AUC of precision-recall curve. & More informative than AUROC under class imbalance; penalises false positives in the SENSITIVE class. \\
Cohen's $\kappa$ & $(P_{\text{obs}} - P_{\text{chance}}) / (1 - P_{\text{chance}})$ & Label agreement beyond chance; penalises confidently incorrect predictions; scale-invariant. \\
Spearman $\rho$ & Rank correlation between confidence and ground-truth z-score. & Assesses whether confidence is monotonically related to effect size. \\
\bottomrule
\end{tabularx}
\end{table}

\subsubsection{Faithfulness Evaluation}
\label{subsec:faithfulness}

Faithfulness is evaluated using an LLM-as-judge protocol applied to all 31 gold-standard cases. For each decisive agent in Round 1, a judge LLM (Gemini 3.1 Flash Lite) receives the agent's structured \texttt{key\_findings} and free-text reasoning. The judge decomposes the reasoning into atomic factual claims and classifies each as \textit{supported} (directly traceable to a retrieved key\_finding), \textit{inferred} (biomedical domain knowledge derived from supported facts), or \textit{unsupported} (asserts a specific fact not present in the retrieved evidence). Faithfulness $= $ supported $/ $ (supported $+$ unsupported); inferred claims are excluded from the denominator.

\subsection{Ablation Study Design}

Fourteen ablation conditions isolate the contribution of each architectural component. Each ablation modifies exactly one component while keeping all others intact.

\begin{table}[H]
\centering
\caption{All fourteen ablation conditions.}
\label{tab:ablation-design}
\begin{tabularx}{\textwidth}{lX}
\toprule
\textbf{Condition} & \textbf{Modification and rationale} \\
\midrule
\texttt{full\_system} & No modification (baseline). \\
\texttt{no\_debate} & Replace DebateEngine with majority-vote aggregator; no Round 2 critique. Tests whether two-round debate adds value over simple majority voting. \\
\texttt{no\_axioms} & Replace AxiomResolver with confidence-only resolver (highest GCS wins). Tests whether axiom hierarchy adds value over pure confidence ordering. \\
\texttt{random\_axiom\_order} & Axiom hierarchy randomly shuffled at runtime. Tests whether the specific tier ordering matters. \\
\texttt{monolithic\_llm} & Single LLM receives all evidence as natural-language prose; no structured MCP JSON, no multi-agent debate. Tests whether multi-agent architecture and MCP grounding add value over a single LLM with all data. \\
\texttt{no\_mcp} & Agents receive evidence as natural-language prose instead of structured MCP outputs; multi-agent debate retained. Tests whether MCP grounding adds value independently of the debate protocol. \\
\texttt{only\_genomics} & Only T1 Genomics agent active. Standalone contribution of structural mutation evidence. \\
\texttt{only\_transcriptomics} & Only T2 Transcriptomics agent active. Standalone contribution of mRNA expression evidence. \\
\texttt{only\_pathway} & Only T3 Pathway agent active. Standalone contribution of pathway topology evidence. \\
\texttt{only\_pharmacology} & Only T4 Pharmacology agent active. Standalone contribution of the IC50 pharmacological prior. \\
\texttt{no\_genomics} & Remove T1; three-agent debate. Marginal contribution of structural mutation evidence. \\
\texttt{no\_transcriptomics} & Remove T2; three-agent debate. Marginal contribution of mRNA expression evidence. \\
\texttt{no\_pathway} & Remove T3; three-agent debate. Marginal contribution of pathway topology evidence. \\
\texttt{no\_pharmacology} & Remove T4; three-agent debate. Marginal contribution of the IC50 pharmacological prior. \\
\bottomrule
\end{tabularx}
\end{table}

\section{Results}
\label{sec:results}

The full DMAS framework was evaluated on 31 literature-validated gold-standard (cell line, drug) pairs. Evaluation uses the system's deployed configuration: four specialist LLM agents (T1--T4), two-round peer debate with mechanistic float peer scoring, and the endorsement-primary AxiomResolver. All 31 cases received a decisive verdict; coverage is 100\%. The framework achieved 87.1\% accuracy (27/31 correct), AUROC = 0.819, AUPRC = 0.766, Cohen's $\kappa$ = 0.733 (substantial agreement, Landis--Koch scale), sensitivity (TPR) = 1.000, and specificity (TNR) = 0.714.

\subsection{Comparison to Baselines}
\label{sec:baselines}

We compare the full DMAS system against two classes of baseline: (i) a traditional machine-learning baseline trained on pharmacogenomic data, and (ii) a monolithic LLM baseline that receives all evidence but uses no multi-agent structure or debate.

\paragraph{Random Forest (ML baseline).}
We trained a Random Forest classifier on 108,712 (cell line, drug) pairs from the CTRPv2 pharmacogenomic screen using three feature groups per case: target gene mutation (binary), target gene expression z-score (CCLE), and binary driver mutation flags for 21 cancer genes. Training labels were derived from CTRPv2 z-scores (z $<$ -0.5 = SENSITIVE, z $>$ 0.5 = RESISTANT). The model was then evaluated on the same 31 gold-standard cases using literature-derived labels (independent of CTRPv2). This setup mirrors the strongest conventional ML approach to this problem.

The RF achieves AUROC = 0.475 and $\kappa$ = $-$0.016 on the 31 gold-standard cases --- below random on AUROC and near-zero agreement on $\kappa$. The failure is systematic, not incidental: binary mutation flags cannot distinguish resistance-conferring variants (e.g., EGFR T790M) from sensitising ones (e.g., EGFR L858R), and the population-level rules learned from CTRPv2 z-scores do not generalise to mechanistically complex cases where the resistance phenotype is published because it is an exception to the population trend.

\paragraph{Monolithic LLM baseline.}
A single LLM (same model as the DMAS agents) receives all omics evidence as natural-language prose --- genomic, transcriptomic, pathway, and pharmacological --- with no structured MCP tool outputs, no specialist agent decomposition, and no debate protocol. It achieves AUROC = 0.656 and $\kappa$ = 0.491. The full DMAS system improves over this baseline by 16.3 AUROC points and 24.2 $\kappa$ points, demonstrating that the multi-agent structure and debate protocol are not cosmetic --- they produce substantively better mechanistic reasoning.

\paragraph{Published multi-omics and multi-agent models.}
Two published systems are architecturally related but address different tasks and cannot be directly compared on our benchmark.
\textit{DRN-CDR}~\cite{saranya2024} trains a deep residual network to predict continuous IC50 values on GDSC (Pearson $r$ = 0.794, RMSE = 0.92). It performs IC50 interpolation within the GDSC training distribution and requires DNA methylation features not in our pipeline; it does not produce binary mechanistic verdicts. Our RF proxy result (AUROC = 0.475) suggests that even a fully fitted model from the same paradigm fails on mechanistically complex out-of-distribution cases.
\textit{EvoMDT}~\cite{liu2026} is a self-evolving multi-agent system for oncology clinical treatment recommendation (MedQA-Cancer accuracy = 0.89). It shares architectural parallels with our system --- specialist agent roles, evidence-grounded reasoning, conflict resolution --- but targets clinical QA rather than pre-clinical mechanistic drug-response prediction. These systems are complementary rather than competing.

\begin{table}[H]
\centering
\caption{Baseline comparison on the 31 gold-standard cases. The RF and monolithic LLM baselines were run on the same cases and labels as the full system. Published models (DRN-CDR, EvoMDT) address different tasks and are included for context only.}
\label{tab:baseline-comparison}
\begin{tabular}{lrrll}
\toprule
\textbf{Method} & \textbf{AUROC} & \textbf{$\kappa$} & \textbf{Task} & \textbf{Notes} \\
\midrule
\textbf{DMAS (ours)} & \textbf{0.819} & \textbf{0.733} & Binary, 31 gold-std & \\
Monolithic LLM (ours) & 0.656 & 0.491 & Binary, 31 gold-std & No agents, no debate \\
RF on CTRPv2 omics (ours) & 0.475 & $-$0.016 & Binary, 31 gold-std & 108K training pairs \\
\midrule
DRN-CDR~\cite{saranya2024} & --- & --- & IC50 regression (GDSC) & Different task; Pearson $r$=0.794 \\
EvoMDT~\cite{liu2026} & --- & --- & Clinical QA & Different task; MedQA-Cancer=0.89 \\
\bottomrule
\end{tabular}
\end{table}

\subsection{Accuracy and Resolution Breakdown}

% TODO: INSERT FIGURE --- Figure 5: Resolution breakdown across 31 cases. Bar chart or pie chart showing CONSENSUS_R1=17, CONSENSUS_R2=8, RESOLVER_TIEBREAK=6.
\begin{figure}[H]
  \centering
  \caption{Resolution breakdown: CONSENSUS\_R1 = 54.8\% (17/31), CONSENSUS\_R2 = 25.8\% (8/31), RESOLVER\_TIEBREAK = 19.4\% (6/31).}
  \label{fig:resolution}
\end{figure}

Of the 31 cases, 54.8\% (17/31) reached consensus in Round 1; 25.8\% (8/31) reached consensus in Round 2 after peer critique; and 19.4\% (6/31) required the AxiomResolver. The R2 verdict revision rate is 19.6\% (11 of 56 agent instances revised their Round 1 verdict in Round 2). The average mechanistic peer score is 0.699.

% TODO: INSERT FIGURE --- Figure 6: Winning agent distribution pie chart. T1 Genomics 87.1% (27/31), T3 Pathway 6.5% (2/31), T2 Transcriptomics 3.2% (1/31), T4 Pharmacology 3.2% (1/31).
\begin{figure}[H]
  \centering
  \caption{Winning agent distribution: T1 Genomics drove 87.1\% of verdicts (27/31), reflecting the strong published mechanistic evidence in the gold-standard case set.}
  \label{fig:winning-agent}
\end{figure}

T1 Genomics drove 87.1\% of verdicts (27/31). The gold-standard cases were specifically selected because their resistance phenotypes are mechanistically published --- the literature that defines the ground truth is precisely the mutation-level, CIViC-annotated evidence that T1 retrieves.

\subsection{Confidence Calibration}

% TODO: INSERT FIGURE --- Figure 7: GCS confidence vs. outcome scatter plot for all 31 cases. X-axis: GCS confidence (0-1). Y-axis: binary correct/wrong. Highlight MCF7+Palbociclib as the highest-confidence wrong prediction at 0.939.
\begin{figure}[H]
  \centering
  \caption{GCS confidence vs.\ outcome for all 31 gold-standard cases. The system makes 4 wrong predictions. The highest-confidence error is MCF7 + Palbociclib (GCS = 0.939, predicted SENSITIVE, true label RESISTANT) due to a cytostatic assay artifact.}
  \label{fig:calibration}
\end{figure}

The system makes 4 wrong predictions across 31 cases. The most striking is \textbf{MCF7 + Palbociclib} (GCS = 0.939, predicted SENSITIVE, true label RESISTANT). Palbociclib is a CDK4/6 inhibitor that arrests MCF7 cells in G1 without killing them. Standard viability assays measure cell number and report apparent sensitivity because growth is inhibited, while cells are biologically resistant --- they survive, resume cycling upon drug withdrawal, and are not killed. This mechanism is undetectable from IC50 z-scores or CIViC mutation evidence alone; it requires assay-type metadata (viability vs.\ cytotoxicity) not present in the current evidence schema.

\subsection{Ablation Study}

Fourteen ablation conditions were conducted on the same 31 gold-standard cases. Each condition removes or replaces exactly one architectural component (Table~\ref{tab:ablation-results}).

\begin{table}[H]
\centering
\caption{Performance of the full system and all 14 ablation conditions on 31 gold-standard cases.}
\label{tab:ablation-results}
\begin{tabular}{lrrrrc}
\toprule
\textbf{Condition} & \textbf{AUROC} & \textbf{AUPRC} & \textbf{Cohen's} $\boldsymbol{\kappa}$ & \textbf{Spearman} $\boldsymbol{\rho}$ & \textbf{Cov.} \\
\midrule
\texttt{full\_system} & 0.819 & 0.766 & 0.733 & 0.551 & 100\% \\
\midrule
\texttt{no\_debate} & 0.784 & 0.799 & 0.668 & --- & 100\% \\
\texttt{no\_axioms} & 0.803 & 0.805 & 0.668 & --- & 100\% \\
\texttt{random\_axiom\_order} & 0.808 & 0.807 & 0.616 & --- & 97\% \\
\texttt{monolithic\_llm} & 0.656 & 0.654 & 0.491 & --- & 97\% \\
\texttt{no\_mcp} & 0.817 & 0.857 & 0.537 & --- & 97\% \\
\midrule
\texttt{only\_genomics} & 0.772 & 0.676 & 0.489 & --- & 84\% \\
\texttt{only\_transcriptomics} & 0.709 & 0.856 & 0.215 & --- & 55\% \\
\texttt{only\_pathway} & 0.673 & 0.633 & 0.266 & --- & 83\% \\
\texttt{only\_pharmacology} & 0.863 & 0.912 & 0.572 & --- & 94\% \\
\midrule
\texttt{no\_genomics} & 0.831 & 0.808 & 0.520 & --- & 97\% \\
\texttt{no\_transcriptomics} & 0.815 & 0.813 & 0.535 & --- & 100\% \\
\texttt{no\_pathway} & 0.824 & 0.836 & 0.604 & --- & 100\% \\
\texttt{no\_pharmacology} & 0.742 & 0.729 & 0.599 & --- & 100\% \\
\bottomrule
\end{tabular}
\end{table}

% TODO: INSERT FIGURE --- Figure 8: Grouped horizontal bar chart of AUROC and kappa for all 14 ablation conditions, with dashed reference lines at full_system values.
\begin{figure}[H]
  \centering
  \caption{AUROC and Cohen's $\kappa$ for all 14 ablation conditions. Dashed lines indicate full system values (AUROC = 0.819, $\kappa$ = 0.733).}
  \label{fig:ablation-chart}
\end{figure}

\textbf{Key ablation findings:}
\begin{itemize}
  \item \textbf{\texttt{no\_debate}:} Coverage remains 100\% with the majority-vote engine. Removing debate reduces AUROC from 0.819 to 0.784 and $\kappa$ from 0.733 to 0.668 --- a 65-point drop in per-case agreement. The debate's value is discriminative quality improvement, not coverage rescue.
  \item \textbf{\texttt{monolithic\_llm}:} The weakest condition (AUROC = 0.656, $\kappa$ = 0.491), confirming that both the multi-agent architecture and the MCP grounding layer contribute independently.
  \item \textbf{\texttt{only\_pharmacology}:} Highest single-agent AUROC (0.863) but $\kappa$ = 0.572, failing on 5 mechanistic edge cases involving synthetic lethality, cytostatic artifacts, and bypass blindness.
  \item \textbf{\texttt{only\_transcriptomics}:} Lowest coverage (55\%), confirming that expression data alone cannot resolve cases where the drug target is wild-type.
  \item \textbf{\texttt{no\_axioms}:} AUROC = 0.803, $\kappa$ = 0.668 --- both strictly lower than full system. The axiom hierarchy demonstrably improves both metrics.
  \item \textbf{\texttt{no\_pharmacology}:} AUROC = 0.742, $\kappa$ = 0.599. The largest single-agent marginal degradation.
\end{itemize}

\subsection{Representative Debate Traces}

Three cases illustrate the qualitatively distinct debate outcomes:

% TODO: INSERT FIGURE --- Figure 9: Clean R1 consensus trace. A case where T1 Genomics finds a CIViC-annotated driver mutation and all agents or the majority agree in Round 1 without needing Round 2.
\begin{figure}[H]
  \centering
  \caption{Clean R1 consensus driven by T1 structural evidence.}
  \label{fig:trace1}
\end{figure}

% TODO: INSERT FIGURE --- Figure 10: Minority-agent-wins trace. A case (e.g. MDA-MB-436+Olaparib or 8505C+PLX-4720) where one agent holds a minority verdict in R1 but wins via AxiomChallenge or resolver after debate.
\begin{figure}[H]
  \centering
  \caption{Minority agent wins after debate: AxiomChallenge forces R2 escalation; resolver selects the minority agent via axiom priority.}
  \label{fig:trace2}
\end{figure}

% TODO: INSERT FIGURE --- Figure 11: T4 pharmacology as sole decisive agent in a case where T1-T3 are all UNCERTAIN.
\begin{figure}[H]
  \centering
  \caption{T4 pharmacology as sole decisive agent: T1--T3 all return UNCERTAIN; T4 IC50 label is the only signal.}
  \label{fig:trace3}
\end{figure}

\subsection{Reasoning Quality Evaluation}
\label{subsec:rq-results}

Faithfulness was evaluated using the LLM-as-judge protocol (Section~\ref{subsec:faithfulness}) on all 31 gold-standard cases.

Across 194 atomic claims from 47 decisive agent instances, the system achieves a claim-level faithfulness of \textbf{86.6\%} (168/194 supported). The agent-level hallucination rate --- fraction of agents with at least one unsupported claim --- is 38.3\% (18/47). Per-agent faithfulness: T1 Genomics = 0.929, T4 Pharmacology = 0.925, T2 Transcriptomics = 0.812, T3 Pathway = 0.806.

The primary failure mode is \textbf{cross-tier leakage}: agents cite mechanistically correct facts not present in the data returned by their assigned MCP tools, but present in the model's parametric knowledge. This is distinct from hallucination --- the facts are not fabricated, but they are not grounded in the retrieved evidence. For example, a T3 Pathway agent may correctly state that a gene participates in the MAPK/ERK cascade, but if that information was not returned by \texttt{find\_pathways\_for\_gene} for that specific case, the claim is unsupported.

% TODO: INSERT FIGURE --- Figure 12: Per-agent faithfulness bar chart (T1=0.929, T4=0.925, T2=0.812, T3=0.806) and winning agent distribution (T1=87.1%, T3=6.5%, T2=3.2%, T4=3.2%).
\begin{figure}[H]
  \centering
  \caption{Per-agent faithfulness scores and winning agent distribution across 31 gold-standard cases.}
  \label{fig:faithfulness}
\end{figure}

\section{Discussion}
\label{sec:discussion}

The full DMAS framework achieved 87.1\% accuracy (27/31 correct) on the 31 gold-standard cases, with AUROC = 0.819, AUPRC = 0.766, Cohen's $\kappa$ = 0.733, coverage = 100\%, sensitivity (TPR) = 1.000, and specificity (TNR) = 0.714. A $\kappa$ of 0.733 represents substantial agreement on the Landis--Koch scale (0.61--0.80). Sensitivity of 1.000 --- zero false negatives --- is the clinically most important result: the system never falsely classifies a sensitive cell line as resistant, which in a clinical context would mean withholding an effective therapy.

\textbf{Key finding:} The DMAS achieves competitive performance (AUROC = 0.819, $\kappa$ = 0.733) without supervised training on labelled drug-resistance data. Performance is limited by database completeness and assay-type metadata gaps, not by the debate architecture or agent reasoning quality.

\subsection{The Debate Protocol: Quality as the Primary Contribution}

The most important result from the ablation study concerns the effect of removing the debate engine entirely (\texttt{no\_debate}). Contrary to earlier analysis, coverage does \emph{not} collapse: the majority-vote engine achieves 100\% coverage. This finding reframes the debate protocol's value. The two-round peer debate is not necessary for the system to produce a verdict on every case --- it is necessary for the system to produce \textit{correct} verdicts on hard cases.

Removing debate reduces AUROC from 0.819 to 0.784 and $\kappa$ from 0.733 to 0.668, a 65-point drop in per-case agreement beyond chance. The debate is architecturally valuable not because it rescues coverage but because it suppresses overconfident single-agent verdicts in mechanistically ambiguous cases. When agents observe in Round 2 that one or more peers have committed to a verdict with stated reasoning and cited evidence, the peer critique process creates an evidence-based context that prompts agents to re-evaluate. This is not sycophantic convergence; it reflects a genuine evidence-backed update that multi-tier integration provides: knowing that the pathway agent found no bypass routes (negative evidence) combined with the pharmacology agent finding a low z-score IC50 (positive evidence) gives more information than either agent alone.

\textbf{Key finding:} The debate engine's primary value is discriminative quality improvement ($\kappa$ +65 points, AUROC +3.5 points over majority-vote), not coverage rescue.

\subsection{Comparison to Existing Methods}

The monolithic LLM baseline (\texttt{monolithic\_llm}) --- a single LLM receiving all evidence as natural-language prose with no structured MCP grounding and no multi-agent debate --- achieves AUROC = 0.656 and $\kappa$ = 0.491. The full DMAS system improves over this baseline by 16.3 AUROC points and 24.2 $\kappa$ points. The structured debate and MCP grounding each contribute independently: \texttt{no\_mcp} (debate retained) gives AUROC = 0.817 and $\kappa$ = 0.537, while \texttt{no\_debate} (MCP retained) gives AUROC = 0.784 and $\kappa$ = 0.668. The largest marginal contribution is from the debate protocol ($\Delta\kappa = +0.196$ over monolithic LLM) with MCP grounding providing an additional $\Delta\kappa = +0.131$.

\subsection{Agent Role Distribution and Tier Hierarchy}

T1 Genomics drove 27 of 31 verdicts (87.1\%). This reflects the case selection methodology: the 31 gold-standard cases were specifically selected because their resistance phenotypes are established in peer-reviewed publications, and the published mechanistic evidence is precisely the mutation-level, CIViC-annotated, drug-specific evidence that T1 retrieves. In a clinical setting where a novel drug lacks a published mechanistic literature, T4 pharmacological IC50 data and T3 pathway topology would become proportionally more decisive.

The T4 Pharmacology Agent exhibits a paradox: as a standalone condition (\texttt{only\_pharmacology}), T4 achieves the highest AUROC of any single-agent condition (0.863) but relatively modest $\kappa$ (0.572). This reflects T4's strength in ranking cases by IC50 z-score (AUROC measures rank quality) combined with its failure on mechanistic edge cases involving synthetic lethality, cytostatic growth arrest, and bypass blindness --- situations where the IC50-derived label is technically correct but mechanistically misleading.

The T3 Pathway Agent represents the most complex tradeoff. The \texttt{no\_pathway} ablation shows reduced $\kappa$ (0.604 versus 0.733) but comparable AUROC (0.824 versus 0.819). T3 adds categorical correctness ($\kappa$) while its occasional bypass detection misfires can introduce confident incorrect calls. Whether T3 is net positive depends on the clinical use case: ranking potential responders favours including T3, whereas minimising confidently incorrect calls favours weighting T3 lower.

\subsection{Confidence Calibration and High-Confidence Errors}

The GCS confidence formula was designed to compute conservative scores by weighting evidence tier and signal strength in addition to raw biomarker citation rate. The empirical calibration results are mixed.

The most concerning case is \textbf{MCF7 + Palbociclib} (GCS = 0.939, predicted SENSITIVE, true label RESISTANT). The system correctly identifies CDK4/6 pathway activity from CIViC annotations and pathway topology, producing a high-confidence SENSITIVE prediction. However, Palbociclib is cytostatic --- it arrests MCF7 cells in G1 without killing them, and viability assays report apparent sensitivity. This error cannot be resolved by adding more databases; it requires assay-type metadata (viability vs.\ cytotoxicity classification) at the evidence layer.

A formally calibrated confidence score would discount GCS when the winning verdict rests on a single decisive agent without corroborating cross-tier evidence --- a ``singleton penalty'' rule requiring a sole decisive agent to exceed GCS $\geq 0.80$ before producing a decisive output.

% TODO: INSERT FIGURE --- Confidence calibration plot: GCS confidence on x-axis vs. outcome on y-axis for all 31 cases, highlighting the 4 wrong predictions.
\begin{figure}[H]
  \centering
  \caption{Confidence calibration plot. GCS confidence vs.\ outcome for all 31 cases. Wrong predictions are highlighted, with MCF7 + Palbociclib as the highest-confidence error (GCS = 0.939).}
  \label{fig:conf-calibration}
\end{figure}

\subsection{What the Ablation Tells Us About Architecture}

\subsubsection{The debate protocol: quality, not coverage}

Removing the debate protocol (\texttt{no\_debate}) achieves 100\% coverage with the majority-vote engine. The debate's value is in improving the quality of decisions: AUROC increases from 0.784 to 0.819 and $\kappa$ increases from 0.668 to 0.733 when debate is enabled. The mechanism is the AxiomChallenge and peer-pressure revision system: agents that would commit to a wrong verdict in isolation can be corrected when a higher-tier peer challenges them with modality-specific evidence. The 19.6\% R2 revision rate confirms that this mechanism is active.

\subsubsection{The axiom hierarchy: consistent improvement, not paradox}

Replacing the endorsement-primary axiom resolver with pure confidence-weighted selection (\texttt{no\_axioms}) achieves AUROC = 0.803 and $\kappa$ = 0.668 --- both strictly lower than the full system (0.819, 0.733). The axiom hierarchy demonstrably improves both metrics. The \texttt{random\_axiom\_order} condition ($\kappa$ = 0.616, coverage = 97\%) further confirms that the specific tier ordering matters, not merely the presence of any hierarchical constraint.

\subsubsection{The no pathway tension}

The \texttt{no\_pathway} ablation shows comparable AUROC (0.824 versus 0.819) but reduced $\kappa$ (0.604 versus 0.733). T3 improves categorical correctness when present but its occasional bypass detection misfires limit ranking gains. The appropriate response is not to remove T3, but to calibrate its bypass-detection threshold more conservatively --- for example, requiring confirmation from both T3 and at least one other agent before allowing T3 to drive a decisive verdict.

\textbf{Architectural implication:} The ablation results together suggest that the main architectural weakness is single-agent verdict dominance when other agents abstain. A ``confidence floor for singleton verdicts'' rule requiring a sole decisive agent to exceed GCS $\geq 0.80$ before producing a decisive output would likely reduce confidently incorrect cases without substantially degrading coverage.

\section{Limitations}
\label{sec:limitations}

\begin{itemize}
  \item \textbf{Database completeness:} All four MCP databases have coverage gaps for novel or poorly annotated drugs. Resistance mechanisms not represented in CIViC, KEGG, CTRPv2, or CCLE cannot be captured by any agent. This is the primary source of the 4 incorrect predictions out of 31 and is external to the system architecture.

  \item \textbf{Cytostatic assay artifact:} Viability-based IC50 assays cannot distinguish cytostatic growth arrest from cytotoxic cell death. Drugs like Palbociclib arrest cells without killing them; standard assays report apparent sensitivity while the biological outcome is resistance. The system has no mechanism to detect this without assay-type metadata in the evidence schema.

  \item \textbf{Test-set scale:} At $n = 31$, a single misclassification shifts accuracy by 3.2 percentage points. Conclusions about relative performance should be treated as indicative pending a larger gold-standard evaluation set.

  \item \textbf{Single LLM family:} All agents use Gemini 3.1 Flash Lite through Google Vertex AI. Results may differ with other model families or larger models. The effectiveness of the debate protocol may also be model-dependent.

  \item \textbf{Cross-tier leakage as a faithfulness gap:} The primary faithfulness failure mode is cross-tier leakage (38.3\% of agents have $\geq 1$ unsupported claim). Agents cite mechanistically correct facts from parametric knowledge not returned by the assigned MCP tool call. This is distinct from hallucination --- the facts are not fabricated --- but it represents a grounding gap that undermines the anti-hallucination guarantee of the MCP layer. Mitigation requires a stricter source-citation schema in EvidencePack \texttt{key\_findings}, enforcing that every claim references a specific tool response field.

  \item \textbf{Computational cost and latency:} Each case requires four independent agent calls in Round 1, up to 12 peer-critique calls in Round 2, and multiple MCP server queries per agent. Total latency is approximately 30--90 seconds per case.

  \item \textbf{Binary classification only:} The system predicts SENSITIVE, RESISTANT, or UNCERTAIN. It does not predict continuous IC50 values, dose-response curves, combination synergy, or acquired versus intrinsic resistance.
\end{itemize}

\section{Future Work}
\label{sec:future}

\textbf{Medium-term priorities:}
\begin{itemize}
  \item \textbf{Literature-aware T5 agent:} Add a fifth agent tier that queries PubMed and bioRxiv for drug-specific resistance mechanisms not captured in structured databases. This agent would use a semantic-search MCP server over publication abstracts to target paradoxical-feedback and pharmacokinetic-resistance cases.

  \item \textbf{Assay-type metadata layer:} Integrate assay classification metadata (viability vs.\ cytotoxicity vs.\ clonogenic assay) into the pharmacology evidence schema. This would allow the T4 agent and GCS formula to discount IC50 labels derived from viability-only assays for cytostatic drug classes, directly addressing the MCF7 + Palbociclib failure mode.

  \item \textbf{Multi-model comparison:} Evaluate the debate framework using GPT-4o and Claude Sonnet as alternative underlying LLMs to determine whether the effectiveness of the debate protocol is model-agnostic.
\end{itemize}

\textbf{Long-term priorities:}
\begin{itemize}
  \item \textbf{Drug-combination extension:} Extend the framework to predict synergy or antagonism for two-drug combinations, requiring a modified protocol in which each agent assesses mechanism overlap and interaction effects, with new MCP data sources such as SynergxDB and DrugCombDB.
\end{itemize}

\section{Conclusion}
\label{sec:conclusion}

This study presents a Decentralised Multi-Agent System (DMAS) for cancer drug resistance prediction that achieves strong discrimination (AUROC = 0.819, Cohen's $\kappa$ = 0.733 on 31 literature-validated gold-standard cases) without supervised training on labelled resistance data. The system integrates four specialist LLM agents, each linked to a dedicated MCP server, through a structured two-round debate protocol and an AxiomResolver.

The crucial finding is that the debate protocol's value is \textit{discriminative quality}, not coverage. Both the full system and the no-debate majority-vote condition achieve 100\% coverage on the 31 gold-standard cases. Removing the debate protocol reduces $\kappa$ by 65 points (0.668 vs.\ 0.733) and AUROC by 3.5 points (0.784 vs.\ 0.819), demonstrating that debate's architectural contribution is suppressing overconfident single-agent verdicts on mechanistically ambiguous cases, not rescuing the system from abstention.

The system achieves a sensitivity (TPR) of 1.000 --- zero false negatives across all 31 cases --- which is the clinically most important result. No truly sensitive cell line was classified as resistant. Claim-level faithfulness across 194 atomic claims is 86.6\%, with cross-tier leakage identified as a novel failure mode distinct from hallucination: agents cite mechanistically correct facts from parametric knowledge rather than from retrieved MCP data, representing a grounding gap rather than fabrication.

In 4 of 31 wrong predictions, the correct answer either required knowledge absent from CIViC, KEGG, CTRPv2, and CCLE, or required assay-type classification to distinguish cytostatic from cytotoxic drug effects. This finding identifies specific classes of resistance mechanisms requiring new data sources rather than architectural improvements.

The primary contributions include: a debate-driven quality improvement mechanism, an MCP-based structural anti-hallucination layer, a tier-stratified resolution protocol, and a systematic fourteen-condition ablation study that isolates the independent contribution of each architectural component. The AxiomResolver combines peer endorsement scores, biological tier priority, and GCS confidence to output a final verdict with a full debate trace and dissent log. This work provides a scalable, auditable, and biologically faithful architecture for personalized cancer therapy decision support.

\bibliographystyle{plainnat}
\begin{thebibliography}{99}

\bibitem{li2026} Li A, Zhang J, Li L, et al. M3MAD-Bench: Are multi-agent debates really effective across domains and modalities? 2026.

\bibitem{elsayed2026} El-Sayed AF, Darwish SM, Labib NM, et al. Generative deep learning in bioinformatics: a microservices-driven paradigm for big data environment. \textit{J Big Data}, 13, 2026.

\bibitem{galitsky2024} Galitsky BA. LLM-based personalized recommendations in health. \textit{Preprints} 2024, 2024021709, 2024.

\bibitem{lin2025} Lin C and Kuo CF. Roles and potential of large language models in healthcare. \textit{Biomed J}, 48, 2025.

\bibitem{williams2024} Williams CYK, Miao BY, Kornblith AE, et al. Evaluating LLMs for clinical recommendations in the emergency department. \textit{Nat Commun}, 15, 2024.

\bibitem{mendez2019} Mendez D, Gaulton A, Bento AP, et al. ChEMBL: towards direct deposition of bioassay data. \textit{Nucleic Acids Res}, 47(D1):D930--D940, 2019.

\bibitem{carli2025} Carli F, Di Chiaro P, Morelli M, et al. Learning and actioning general principles of cancer cell drug sensitivity. \textit{Nat Commun}, 16, 2025.

\bibitem{iorio2016} Iorio F, Knijnenburg TA, Vis DJ, et al. A landscape of pharmacogenomic interactions in cancer. \textit{Cell}, 166(3):740--754, 2016.

\bibitem{widjaja2025} Widjaja F, Chen Z, and Zhou J. BioinfoMCP: A unified platform enabling MCP interfaces in agentic bioinformatics, 2025.

\bibitem{housman2014} Housman G, Byler S, Heerboth S, et al. Drug resistance in cancer: an overview. \textit{Cancers}, 6(3):1769--1792, 2014.

\bibitem{he2025} He H et al. Biology-Instructions: A dataset and benchmark for multi-omics sequence understanding of LLMs. In \textit{Findings of EMNLP 2025}, 2025.

\bibitem{zhao2023} Zhao H, Zhang X, Zhao Q, Li Y, and Wang J. MSDRP: a deep learning model for drug response prediction. \textit{Bioinformatics}, 39(9), 2023.

\bibitem{chen2024} Chen HO, Cui YC, Lin PC, and Chiang JH. Multi-omics model with latent alignment for drug response prediction. \textit{J Pers Med}, 14, 2024.

\bibitem{proietti2020} Proietti I, Skroza N, Bernardini N, et al. Mechanisms of acquired BRAF inhibitor resistance in melanoma. \textit{Cancers}, 12(10), 2020.

\bibitem{barretina2012} Barretina J, Caponigro G, Stransky N, et al. The cancer cell line encyclopedia enables predictive modelling of anticancer drug sensitivity. \textit{Nature}, 483:603--607, 2012.

\bibitem{rotow2017} Rotow J and Bivona TG. Understanding and targeting resistance mechanisms in NSCLC. \textit{Nat Rev Cancer}, 17(11):637--658, 2017.

\bibitem{spieser2026} Spieser J, Balapour A, Meller J, Patra KC, and Shamsaei B. A review of multi-agent AI systems for biological and clinical data analysis. \textit{Methods Protoc}, 9(2), 2026.

\bibitem{zhou2025streamline} Zhou J, Jiang J, Han Z, Wang Z, and Gao X. Streamline automated biomedical discoveries with agentic bioinformatics. \textit{Brief Bioinform}, 26(5), 2025.

\bibitem{vasilev2025} Vasilev K, Misrahi A, Jain E, et al. MTBBench: A multimodal sequential clinical decision-making benchmark in oncology, 2025.

\bibitem{saranya2024} Saranya KR and Vimina ER. DRN-CDR: A cancer drug response prediction model using multi-omics. \textit{Comput Biol Chem}, 112, 2024.

\bibitem{lorenzo2026} Lorenzo L, Montana-Mendez M, Figueiras S, Boubeta M, and Bernardo-Castineira C. Evaluation of OncoTimia for supporting tumour boards, 2026.

\bibitem{parca2025} Parca L, Petrizzelli M, Conte F, et al. CellHit: predicting patient drug responses via cancer cell lines. \textit{Nucleic Acids Res}, 2025.

\bibitem{kaesberg2025} Kaesberg LB, Becker J, Wahle JP, Ruas T, and Gipp B. Voting or consensus? Decision-making in multi-agent debate. In \textit{Findings of ACL 2025}, Vienna, 2025.

\bibitem{flotho2026} Flotho M, Diks IF, Flotho P, et al. MCPmed: a call for model context protocol-enabled bioinformatics web services. \textit{Brief Bioinform}, 27(1), 2026.

\bibitem{griffith2017} Griffith M, Spies N, Krysiak K, et al. CIViC is a community knowledgebase for expert crowdsourcing clinical interpretation of variants in cancer. \textit{Nature Genetics}, 49:170--174, 2017.

\bibitem{kanehisa2000} Kanehisa M and Goto S. KEGG: Kyoto encyclopedia of genes and genomes. \textit{Nucleic Acids Res}, 28(1):27--30, 2000.

\bibitem{turpin2023} Turpin M, Michael J, Perez E, and Bowman SR. Language models don't always say what they think. In \textit{NeurIPS}, 2023.

\bibitem{zack2025} Zack M, Stupichev DN, Moore AJ, et al. AI and multi-omics in pharmacogenomics. \textit{Mayo Clin Proc Digit Health}, 3(3), 2025.

\bibitem{rees2016} Rees MG, Seashore-Ludlow B, Cheah JH, et al. Correlating chemical sensitivity and basal gene expression reveals mechanism of action. \textit{Nature Chemical Biology}, 12:109--116, 2016.

\bibitem{codella2025} Codella N, Preston S, Qiu H, et al. Healthcare agent orchestrator (HAO) for patient summarization in molecular tumor boards, 2025.

\bibitem{mehandru2025} Mehandru N, Hall AK, Melnichenko O, et al. BioAgents: Bridging the gap in bioinformatics with multi-agent systems. \textit{Sci Rep}, 15, 2025.

\bibitem{du2024} Du P, Fan R, Zhang N, Wu C, and Zhang Y. Advances in integrated multi-omics analysis for drug-target identification. \textit{Biomolecules}, 14(6), 2024.

\bibitem{liu2026evomdt} Liu Q, Hu Z, Huang T, et al. EvoMDT: a self-evolving multi-agent system for clinical decision-making in multi-cancer. \textit{npj Digit Med}, 9, 2026.

\bibitem{vavekanand2026} Vavekanand R. A comprehensive review of multimodal LLMs for medical imaging and omics data. \textit{Arch Comput Methods Eng}, 2026.

\bibitem{zhang2025} Zhang R, Wang J, Yang S, et al. Multi-agent intelligence for multidisciplinary decision-making in gastrointestinal oncology, 2025.

\bibitem{ali2025} Ali S, Qadri YA, Ahmad K, et al. Large language models in genomics. \textit{Bioengineering (Basel)}, 12(5), 2025.

\bibitem{qu2025} Qu S, Ding N, Xie L, et al. Automating exploratory multiomics research via language models, 2025.

\bibitem{zhou2025scoping} Zhou S, Xu Z, Zhang M, et al. Large language models for disease diagnosis: a scoping review. \textit{npj Artif Intell}, 1, 2025.

\bibitem{almutiri2023} Almutiri T, Alomar K, and Alganmi N. Predicting drug response on multi-omics data using Bayesian ridge regression with deep forest. \textit{Int J Adv Comput Sci Appl}, 14, 2023.

\bibitem{liu2024groupdebate} Liu T, Wang X, Huang W, et al. GroupDebate: Enhancing the efficiency of multi-agent debate using group discussion, 2024.

\bibitem{yang2025rise} Yang T, Xiao Y, Bao Z, Hao J, and Peng J. The rise and potential of LLM agents in bioinformatics and biomedicine. \textit{Brief Bioinform}, 26(6), 2025.

\bibitem{martinek2025} Martinek V, Gariboldi A, Tzimotoudis D, et al. Agentomics-ML: Autonomous ML experimentation agent for genomic and transcriptomic data, 2025.

\bibitem{jiang2025} Jiang W, Ye W, Tan X, et al. Network-based multi-omics integrative analysis in drug discovery. \textit{BioData Mining}, 18, 2025.

\bibitem{han2025} Han X, Gao X, Qu X, and Yu Z. Multi-agent medical decision consensus matrix system for oncology MDT consultations, 2025.

\bibitem{liu2025multiomics} Liu X, Wei K, Lin E, et al. Multi-omics advances in understanding cancer drug resistance. \textit{Biomed Technol}, 12, 2025.

\bibitem{meng2024} Meng X, Yan X, Zhang K, et al. The application of large language models in medicine: a scoping review. \textit{iScience}, 27(5), 2024.

\bibitem{wu2026} Wu X, Sun M, Li L, Wang J, and Wan S. Leveraging AI and LLMs for cancer immunotherapy. \textit{Advanced Science}, 2026.

\bibitem{yang2025promptbio} Yang X, Shashidhar KC, Zhang M, et al. PromptBio: A multi-agent AI platform for bioinformatics data analysis, 2025.

\bibitem{liu2023snf} Liu XY and Mei XY. Prediction of drug sensitivity based on multi-omics data using deep learning and SNF approaches. \textit{Front Bioeng Biotechnol}, 11, 2023.

\bibitem{cui2025} Cui Y, Fu H, Zhang H, Wang L, and Zuo C. Free-MAD: Consensus-free multi-agent debate, 2025.

\bibitem{kim2025} Kim Y, Jeong H, Chen S, et al. Medical hallucination in foundation models and their impact on healthcare, 2025.

\bibitem{liu2026heterogeneous} Liu Y, Yao H, Liu W, He A, and Zhang Y. Heterogeneous consensus-progressive reasoning for efficient multi-agent debate, 2026.

\bibitem{mao2025} Mao Y, Shangguan D, Huang Q, et al. Emerging AI-driven precision therapies in tumor drug resistance. \textit{Mol Cancer}, 24(1), 2025.

\bibitem{su2025} Su Y, Liu J, Li J, et al. Applications of AI and LLMs in cancer immunotherapy. \textit{Current Proteomics}, 22(5), 2025.

\bibitem{lai2025} Lai YJ, Wang LJ, Yasaka TM, et al. Inferring drug-gene relationships in cancer using literature-augmented LLMs. \textit{Cancer Res Commun}, 5(4):706--718, 2025.

\bibitem{hein2025} Hein ZM, Guruparan D, Okunsai B, et al. AI and machine learning in biology: From genes to proteins. \textit{Biology}, 14, 2025.

\end{thebibliography}

\end{document}
